\section{Moving witnesses to a narrower box}
\label{sec:translations}

A recursive child will receive a box centered at a pivot $k$ near the
middle of its parent box. To transfer a parent's obligations to this
child, we translate a witness by one scalar and clip it at a boundary.
This section establishes the exact support and claims of the
translated vector. It is the geometric part of the argument and uses
only monotonicity, common-shift invariance, and the witness definition.

Consider an input $(A,D,a,b)$ with dyadic $D\ge4096$. In addition
to $u=A+D$ and $P=D/32$, fix
\begin{equation}
 c=A+D/2,\qquad M=D/8,\qquad q=D/128,\qquad
 p=q/32=D/4096.
 \label{eq:pivot-parameters}
\end{equation}
The vector $c$ is the center, $M$ bounds the pivot's displacement,
$q$ is the narrow child width, and $p$ will be one pivot step. All
these quantities, as well as $q/2$ and $D/4$, are integral. We
consider integer pivots $c-M\le k\le c+M$.

For a relevant Max witness $x$, let
\begin{equation}
 s_x(k)=\max_{i\in V}(x_i-k_i),\qquad
 g_i^x(k)=s_x(k)-(x_i-k_i).
 \label{eq:max-gaps}
\end{equation}
We call $s_x(k)$ its \emph{peak} and $g_i^x(k)$ its gap at $i$.
For a relevant Min witness $y$ the dual definitions are
\begin{equation}
 s_y(k)=\max_{i\in V}(k_i-y_i),\qquad
 g_i^y(k)=s_y(k)-(k_i-y_i).
 \label{eq:min-gaps}
\end{equation}
Every gap is a nonnegative integer. The following lemma treats the
two widths that the recursion will use.

\begin{lemma}[Translated supports and claims]
\label{lem:translation}
Let $(A,D,a,b)$ be an input with $D\ge4096$, let
$c-M\le k\le c+M$ be integral, and let $r\in\{q,D/2\}$.
For any relevant Max witness $x$, the vector
\begin{equation}
 x_i'=\max\{k_i-r/2,\ x_i-s_x(k)+r/2\}
 \label{eq:max-translation}
\end{equation}
is an integer subsolution in $[k-r/2,k+r/2]$. Its support is
\[
                  S_{x'}=\{i:g_i^x(k)<r\}\subseteq S_x.
\]
For any positive new Max masses $a'$ with $a'(S_{x'})\le1$,
it is a Max witness and claims every $i$ with $g_i^x(k)\le r/32$.

For any relevant Min witness $y$, the vector
\begin{equation}
 y_i'=\min\{k_i+r/2,\ y_i+s_y(k)-r/2\}
 \label{eq:min-translation}
\end{equation}
is an integer supersolution in the same box, with support
\[
                  S_{y'}=\{i:g_i^y(k)<r\}\subseteq S_y.
\]
For any positive new Min masses $b'$ with $b'(S_{y'})\le1$,
it is a Min witness and claims every $i$ with $g_i^y(k)\le r/32$.
In particular, the corresponding unchanged masses $a$ and $b$ always
make these translated vectors witnesses.
\end{lemma}

\begin{proof}
The support assertion requires knowing which coordinates could survive
the clipping. Relevance supplies that information. A Max claim at
$j$ gives $x_j\ge A_j+D-P$, while $k_j\le A_j+D/2+M$;
hence
\begin{equation}
                         s_x(k)\ge D/2-M-P.
 \label{eq:peak-lower}
\end{equation}
If $i\notin S_x$, then $x_i=A_i$ and
$x_i-k_i\le-D/2+M$. Combining the two inequalities gives
\begin{equation}
                  g_i^x(k)\ge D-2M-P=23D/32>D/2.
 \label{eq:outside-gap}
\end{equation}
The untruncated term in Equation~\eqref{eq:max-translation} is
\[
                   k_i+r/2-g_i^x(k).
\]
It lies below or at the new upper boundary. Taking its maximum with
the lower boundary puts $x'$ inside the new box, with
support exactly where $g_i^x(k)<r$. Since $r\le D/2$,
Equation~\eqref{eq:outside-gap} shows that this support lies in
$S_x$. Integrality follows from the integral half-widths and peaks.

Write $t=-s_x(k)+r/2$. Globally $x'\ge x+t$; on its new
support, $x_i'=x_i+t$ and the parent's subsolution inequality
applies. Equation~\eqref{eq:operator-properties} therefore yields
\[
                 F_i(x')\ge F_i(x+t)=F_i(x)+t
                                \ge x_i+t=x_i'.
\]
Thus $x'$ is a subsolution. Its witness condition is precisely the
stated mass bound. If $g_i^x(k)\le r/32$, its untruncated entry
is at least $k_i+r/2-r/32$, the new Max claim threshold; taking a
maximum with the lower boundary cannot destroy that claim.

For Min, relevance likewise gives $s_y(k)\ge D/2-M-P$.
Outside $S_y$, one has $y_i=A_i+D$, so
$k_i-y_i\le-D/2+M$. The gap is again at least $23D/32>D/2$.
Now the translated term before clipping is
\[
                   k_i-r/2+g_i^y(k).
\]
It is above or at the lower boundary, and taking its minimum with
the upper boundary leaves support exactly at gaps $<r$. The support
is contained in $S_y$ by the same outside-support estimate. Set
$t=s_y(k)-r/2$. The vector inequality $y'\le y+t$ and equality
$y_i'=y_i+t$ on the new support imply
\[
                 F_i(y')\le F_i(y+t)=F_i(y)+t
                                \le y_i+t=y_i'.
\]
This is the required supersolution inequality. The mass condition
makes $y'$ a witness, and a gap at most $r/32$ puts it at most
$r/32$ above the new lower boundary, giving the asserted claim.
Finally, support containment and positivity imply
$a(S_{x'})\le a(S_x)\le1$ and
$b(S_{y'})\le b(S_y)\le1$ for unchanged masses.
\end{proof}

The two thresholds serve different purposes: gaps strictly below $r$
describe the surviving support, whereas gaps at most $r/32$ give
claims.
